\documentclass[letterpaper,10pt,oneside]{article} 
\usepackage[margin=2.54cm]{geometry}
\usepackage[spanish]{babel}
\newcommand\authormark[1]{\textsuperscript{#1}}
\usepackage[colorlinks=true,bookmarks=true,linkcolor=black,citecolor=black,urlcolor=black]{hyperref}
\usepackage{times} 
\usepackage[table]{xcolor}
\usepackage{booktabs}
\usepackage{graphicx}
\usepackage{titlesec}
\usepackage{multicol}
\usepackage{float}
\usepackage{multicol}
\usepackage{multirow}

\titleformat{\section}
{\normalfont\fontsize{12pt}{14pt}\bfseries} % Fuente normal, tamaño 12pt, interlineado 14pt, en negrita
{\thesection}{1em}{} % Formato del número y separación

\begin{document}
	
	\begin{center}
		\fontsize{14}{16}\selectfont\bfseries
		1° Reporte de seguimiento: 
		Simulación de sistemas complejos mediante autómatas celulares neuronales
	\end{center}
	
	\begin{center}
		\fontsize{14}{16}\selectfont\bfseries
		\textit{Trabajo Terminal No. 2026-A114}
	\end{center}
	
	\begin{center}
		\fontsize{12}{16}\selectfont
		\textit{Alumno: \authormark{*}Angeles López Erick Jesse\authormark{1}} \\
		\textit{Director: Juárez Martínez Genaro\authormark{1}} \\
		\textit{\authormark{*}e-mail: eangelesl1700@alumno.ipn.mx}
	\end{center}

	\section{Actividades desarrolladas}
	
\begin{table}[htbp]
	\centering
	\begin{tabular}{|p{0.15\textwidth}|p{0.3\textwidth}|p{0.12\textwidth}|p{0.12\textwidth}|p{0.08\textwidth}|p{0.08\textwidth}|}
		\hline
		\textbf{Actividad} & 
		\textbf{Tarea} &  
		\textbf{Inicio} & 
		\textbf{Finalización} & 
		\textbf{Porcentaje} & 
		\textbf{Prioridad} \\
		\hline
		\multirow{6}{3\textwidth}{\parbox{0.15\textwidth}{\centering Estado del arte}} 
		& Consultar generalidades de autómatas & 20/07/2025 & 15/08/2025 & 100\% & Media \\ \cline{2-6}
		
		& Consultar representaciones morfológicas & 01/08/2025 & 20/08/2025 & 100\% & Alta \\ \cline{2-6}
		& Consultar aplicaciones de sistemas autoorganizados & 10/08/2025 & 20/08/2025 & 100\% & Media \\ \cline{2-6}
		
		& Consultar NCA en mesh & 10/08/2025 & 20/08/2025 & 100\% & Alta \\ \cline{2-6}
		
		& Consultar representaciones adicionales & 20/08/2025 & 30/08/2025 & 100\% & Alta  \\ \cline{2-6}
		
		& Revisión de estado del arte & 30/09/2025 &  & 75\% & Alta  \\ \hline
		
		\multirow{5}{0.15\textwidth}{\parbox{0.15\textwidth}{\centering Marco teórico}} 
		
		& Investigar sobre sistemas complejos & 20/09/2025 & 25/09/2025 & 100\% & Media \\ \cline{2-6}
		
		& Investigar sobre autómatas celulares & 01/09/2025 & 30/09/2025 & 90\% & Alta \\ \cline{2-6}
		
		& Investigar sobre inteligencia artificial y redes neuronales & 10/08/2025 & 25/08/2025 & 100\% & Alta \\ \cline{2-6}
		
		& Investigar y definir autómatas celulares neuronales & 20/09/2025 &  & 80\% & Alta \\ \cline{2-6}
		
		& Revisión de marco teórico & 30/09/2025 &  & 25\% & Alta  \\ \hline
		
		\multirow{5}{0.15\textwidth}{\parbox{0.15\textwidth}{\centering Marco metodológico}} 
		
		& Definir diseño de investigación & &  & 0\% & Alta \\  \cline{2-6}
		
		& Definir planteamiento técnico, metodológico y tecnológico & &  & 0\% & Alta \\  \cline{2-6}
		
		& Definir métodos recolección de datos & &  & 0\% & Alta \\  \cline{2-6}
		
		& Definir métodos de análisis de resultados & &  & 0\% & Alta \\ \cline{2-6}
		
		& Revisión de marco metodológico & &  &  & Alta  \\ \hline
		
		\multirow{3}{0.15\textwidth}{\parbox{0.15\textwidth}{\centering Generalidades de proyecto}} 
		
		& Elaboración del reporte en LaTeX & 20/07/2025 & 20/07/2025 & 100\% & Media \\  \cline{2-6}
		
		& Creación del repositorio del proyecto en GitHub & 20/07/2025 & 20/07/2025 & 100\% & Media \\  \cline{2-6}
		
		& Contacto con sinodales & 20/09/2025 & & 75\% & Alta \\  \hline
		
		
	\end{tabular}
	\caption{Cronograma de actividades}
\end{table}


	
	
	
	\section{Cronograma}
	
	\resizebox{\linewidth}{!}{%
		\begin{tabular}{p{5cm}*{10}{>{\centering\arraybackslash}p{1cm}}}
			\toprule
			\textbf{Actividad} & \textbf{Ago} & \textbf{Sep} & \textbf{Oct} & \textbf{Nov} & \textbf{Dic} & \textbf{Ene} & \textbf{Feb} & \textbf{Mar} & \textbf{Abr} & \textbf{May} \\
			\midrule
			\multicolumn{11}{l}{\textbf{Fase 1: Planificación y Revisión Bibliográfica}} \\
			
			Estado del arte       & \cellcolor{gray!30} & \cellcolor{gray!30} & \cellcolor{gray!30} &  &  &  &  &  &  &   \\
			
			Marco teórico y metodológico &  & \cellcolor{gray!30} & \cellcolor{gray!30} &  &  &  &  &  &  &   \\
			
			\midrule
			
			\multicolumn{11}{l}{\textbf{Fase 2: Diseño e implementación del simulador}} \\
			Arquitectura del simulador &  &  & \cellcolor{gray!30} & \cellcolor{gray!30} &  &  &  &  &  &   \\
			
			Módulos de simulación  &  &  &  & \cellcolor{gray!30} & \cellcolor{gray!30} &  &  &  &  &   \\
			
			Módulos entrenamiento   &  &  &  &  & \cellcolor{gray!30} & \cellcolor{gray!30} &  &  &  &   \\
			
			\midrule
			
			\multicolumn{11}{l}{\textbf{Fase 3: Entrenamiento y Validación}} \\
			Entrenamiento y ajuste   &  &  &  &  &  & \cellcolor{gray!30} & \cellcolor{gray!30} & \cellcolor{gray!30} &   &   \\
			
			Experimentación y análisis  &  &  &  &  &  & \cellcolor{gray!30}  & \cellcolor{gray!30} & \cellcolor{gray!30} & \cellcolor{gray!30} &   \\
			
			Análisis de resultados (global) &  &  &  &  &  & \cellcolor{gray!30} & \cellcolor{gray!30} & \cellcolor{gray!30} & \cellcolor{gray!30} &   \\
			
			\midrule
			
			\multicolumn{11}{l}{\textbf{Fase 4: Documentación y Difusión}} \\
			Documentación técnica  &  &  &  &  &  &  &  & \cellcolor{gray!30}  & \cellcolor{gray!30} & \cellcolor{gray!30}  \\
			
			Artículo académico  &  &  &  &  &  &  &  & \cellcolor{gray!30}  & \cellcolor{gray!30}  & \cellcolor{gray!30}  \\
			
			\bottomrule
		\end{tabular}
	}

	\section{Alumnos y directores}
	
\begin{tabular}{@{} p{0.45\textwidth} p{0.45\textwidth} @{}}
	\textit{Erick Jesse Angeles López}
	&
	\textit{Dr. Genaro Juárez Martínez}
	
	 \\ & \\ & \\
	 & \\ & \\
	 
	 & \\ & \\
	 & \\ & \\
	
	Firma:
	\rule{6cm}{0.4pt} & 
	
	Firma:
	\rule{6cm}{0.4pt} \\
	
	& \\ & \\
	& \\ & \\
	& \\ & \\
	& \\ & \\
\end{tabular}

\end{document}
